Full-law uniqueness identifies \(\mathcal D(P_X)=\{[c_j]:j\in\mathsf S\}\) set-theoretically. For any exact irreducible mixing representative \(C\) of this projective set,
\[
 \boldsymbol\Theta_x^{\mathrm{law}}(P_X)
 =\mathcal V_x(C)
 =\left\{\left(\frac{C_{ij}}{C_{xj}}\right)_{i\in\mathsf O\setminus\{x\}}:j\in J_x(C)\right\},
 \qquad
 \Theta_{y\leftarrow x}^{\mathrm{law}}(P_X)
 =\mathcal R_{xy}(C)
 =\left\{\frac{C_{yj}}{C_{xj}}:j\in J_x(C)\right\}.
\]
These expressions are independent of the representative. Conditional on exact \(C\), the shared-kernel and shared-matching certificate evaluates all frontier values and materializes one exact dense witness per value in worst-case output-optimal \(O(n^3)\) exact arithmetic; this is not a procedure or complexity bound for recovering \(C\) from \(P_X\). The qualitative full-law result does not evaluate \(\mathcal D(P_X)\). Under the heterogeneous theorem's full assumptions---known source count \(n\), normalized full-row-rank pairwise-nonproportional mixing, and mutually independent, nondegenerate, all-non-Gaussian moment-determinate sources---the multi-order handle evaluates the direction set by a terminating pointwise population-oracle procedure. That population-oracle result is incomparable with the accepted lower-order bivariate arrow results because it recovers all projective directions for an arbitrary observed dimension and then maps them to a cyclic response frontier, whereas those results identify a single axis-conditioned bivariate arrow. Under the separated common-order iid assumptions, the declared confidence unions provide simultaneous root-\(N\) inference.